% GENERATED FILE. Edit canonical lessons, then run npm run generate:books.
% canonical-source-hash: fnv1a64:fd781d89047b8623
% canonical-lessons: FR-C443-egout, FR-C443-caniveau, FR-C443-deblayer, FR-C443-saler, FR-C443-deneiger

\chapter{L'hiver en ville}
\label{ch:fr-l-hiver-en-ville}
\hlchaptermodality{443}

\begin{chapteropening}
\textbf{By the end of this chapter:} \emph{I can talk about a town in winter: the sewer, the gutter in the street, clearing the snow, salting the roads and snow-ploughing.}

Complete the last lesson of chapter 443: I can talk about a town in winter: the sewer, the gutter in the street, clearing the snow, salting the roads and snow-ploughing.
\end{chapteropening}

\section[l'égout]{l'égout — the sewer, the drain}

\label{lesson:FR-C443-egout}



\begin{quote}
Before the new one: say the French for the roof gutter, then the French for the covered market.
\end{quote}

\subsection*{You'll want to know: l'égout}
\textbf{l'égout} — \textquotedblleft{}the sewer, the drain\textquotedblright{}.

\begin{quote}
Après l'orage, les égouts sont pleins.
\end{quote}

\begin{quote}
\emph{After the storm the sewers are full.}
\end{quote}

A masculine noun.

\begin{cousinweb}
\emph{égout} is built on \emph{égoutter} 'to let drip, to drain' (EtymDB), from \emph{goutte} 'drop', which goes back to Latin \emph{gutta}.
\end{cousinweb}

\begin{grammarlens}[title={What you've built}]
One more word for telling, explaining and describing at B1.
\end{grammarlens}

\subsection*{Your turn}
\begin{itemize}
\raggedright
  \item \emph{Say it:} \emph{l'égout}
  \item \emph{Say it:} \emph{l'égout}, once more
  \item \emph{Say it:} a sentence of your own with \emph{l'égout}
  \item \emph{Recall:} say the French for the roof gutter, then the French for the covered market, then say \emph{l'égout} again
\end{itemize}

\begin{tcolorbox}[breakable,colback=teal!4,colframe=teal!35!black,title={Before you move on}]
What does l'égout mean? (\textquotedblleft{}The sewer, the drain\textquotedblright{}.) Say it once more.
\end{tcolorbox}
\section[le caniveau]{le caniveau — the gutter (in the street)}

\label{lesson:FR-C443-caniveau}



\begin{quote}
Before the new one: say the French for the covered market, then the French for the sewer.
\end{quote}

\subsection*{You'll want to know: le caniveau}
\textbf{le caniveau} — \textquotedblleft{}the gutter (in the street)\textquotedblright{}.

\begin{quote}
La neige fond et coule dans le caniveau.
\end{quote}

\begin{quote}
\emph{The snow melts and runs into the gutter.}
\end{quote}

At the edge of the road, beside the pavement.

\begin{cousinweb}
The origin of \emph{caniveau} is not settled, and EtymDB does not list it.
\end{cousinweb}

\begin{grammarlens}[title={What you've built}]
One more word for telling, explaining and describing at B1.
\end{grammarlens}

\subsection*{Your turn}
\begin{itemize}
\raggedright
  \item \emph{Say it:} \emph{le caniveau}
  \item \emph{Say it:} \emph{le caniveau}, once more
  \item \emph{Say it:} a sentence of your own with \emph{le caniveau}
  \item \emph{Recall:} say the French for the covered market, then the French for the sewer, then say \emph{le caniveau} again
\end{itemize}

\begin{tcolorbox}[breakable,colback=teal!4,colframe=teal!35!black,title={Before you move on}]
What does le caniveau mean? (\textquotedblleft{}The gutter (in the street)\textquotedblright{}.) Say it once more.
\end{tcolorbox}
\section[déblayer]{déblayer — to clear away (snow, rubble)}

\label{lesson:FR-C443-deblayer}



\begin{quote}
Before the new one: say the French for the sewer, then the French for the gutter.
\end{quote}

\subsection*{You'll want to know: déblayer}
\textbf{déblayer} — \textquotedblleft{}to clear away (snow, rubble)\textquotedblright{}.

\begin{quote}
Les voisins ont déblayé la neige devant l'immeuble.
\end{quote}

\begin{quote}
\emph{The neighbours cleared the snow in front of the building.}
\end{quote}

\emph{Le déblaiement} — the clearing.

\begin{cousinweb}
\emph{déblayer} comes from Old French \emph{desbleer} 'to clear a field of its corn', usually traced to \emph{blé} 'corn'; EtymDB lists the verb without a source.
\end{cousinweb}

\begin{grammarlens}[title={What you've built}]
One more word for telling, explaining and describing at B1.
\end{grammarlens}

\subsection*{Your turn}
\begin{itemize}
\raggedright
  \item \emph{Say it:} \emph{déblayer}
  \item \emph{Say it:} \emph{déblayer}, once more
  \item \emph{Say it:} a sentence of your own with \emph{déblayer}
  \item \emph{Recall:} say the French for the sewer, then the French for the gutter, then say \emph{déblayer} again
\end{itemize}

\begin{tcolorbox}[breakable,colback=teal!4,colframe=teal!35!black,title={Before you move on}]
What does déblayer mean? (\textquotedblleft{}To clear away (snow, rubble)\textquotedblright{}.) Say it once more.
\end{tcolorbox}
\section[saler]{saler — to salt (food, roads)}

\label{lesson:FR-C443-saler}



\begin{quote}
Before the new one: say the French for the gutter, then the French for to clear away.
\end{quote}

\subsection*{You'll want to know: saler}
\textbf{saler} — \textquotedblleft{}to salt (food, roads)\textquotedblright{}.

\begin{quote}
La ville sale les routes quand il gèle.
\end{quote}

\begin{quote}
\emph{The town puts salt on the roads when it freezes.}
\end{quote}

\emph{C'est salé !} — that is pricey (informal).

\begin{cousinweb}
\emph{saler} goes back to Vulgar Latin \emph{salāre}, from Latin \emph{sal} 'salt' (EtymDB).
\end{cousinweb}

\begin{grammarlens}[title={What you've built}]
One more word for telling, explaining and describing at B1.
\end{grammarlens}

\subsection*{Your turn}
\begin{itemize}
\raggedright
  \item \emph{Say it:} \emph{saler}
  \item \emph{Say it:} \emph{saler}, once more
  \item \emph{Say it:} a sentence of your own with \emph{saler}
  \item \emph{Recall:} say the French for the gutter, then the French for to clear away, then say \emph{saler} again
\end{itemize}

\begin{tcolorbox}[breakable,colback=teal!4,colframe=teal!35!black,title={Before you move on}]
What does saler mean? (\textquotedblleft{}To salt (food, roads)\textquotedblright{}.) Say it once more.
\end{tcolorbox}
\section[déneiger]{déneiger — to clear of snow, to plough}

\label{lesson:FR-C443-deneiger}



\begin{quote}
Before the new one: say the French for to clear away, then the French for to salt.
\end{quote}

\subsection*{You'll want to know: déneiger}
\textbf{déneiger} — \textquotedblleft{}to clear of snow, to plough\textquotedblright{}.

\begin{quote}
Le matin, les camions déneigent les grandes rues.
\end{quote}

\begin{quote}
\emph{In the morning the lorries clear the snow from the main streets.}
\end{quote}

Very common in Québec and in the mountains: \emph{le déneigement}.

\begin{cousinweb}
\emph{déneiger} is built from \emph{dé-} and \emph{neige} 'snow', which goes back to Latin \emph{nix, nivis} (Lewis \& Short). EtymDB does not list it.
\end{cousinweb}

\begin{grammarlens}[title={What you've built}]
One more word for telling, explaining and describing at B1.
\end{grammarlens}

\subsection*{Your turn}
\begin{itemize}
\raggedright
  \item \emph{Say it:} \emph{déneiger}
  \item \emph{Say it:} \emph{déneiger}, once more
  \item \emph{Say it:} a sentence of your own with \emph{déneiger}
  \item \emph{Recall:} say the French for to clear away, then the French for to salt, then say \emph{déneiger} again
\end{itemize}

\begin{tcolorbox}[breakable,colback=teal!4,colframe=teal!35!black,title={Before you move on}]
What does déneiger mean? (\textquotedblleft{}To clear of snow, to plough\textquotedblright{}.) Say it once more.
\end{tcolorbox}
